\documentclass[10pt,a4paper]{amsart}
\usepackage[margin=19mm]{geometry}
\usepackage{amsmath,amssymb,mathtools}
\usepackage[scr=rsfs,cal=euler]{mathalfa}
\usepackage{xfrac}
\usepackage[unicode,hidelinks,bookmarks=false]{hyperref}
\newcommand{\ie}{i.e.\ }
\newcommand{\cf}{cf.\ }
\newcommand{\resp}{respectively\ }
\newcommand{\Subsection}{\subsection}
\providecommand{\nobreakdash}{\nobreak\hbox{-}}
\setlength{\emergencystretch}{2em}
\multlinegap0pt
\usepackage{bm}
\usepackage{bbm}
\usepackage{dsfont}
\usepackage{xspace}
\usepackage{cancel}
\usepackage{mathtools}
\usepackage{amsthm}
\makeatletter
\@ifundefined{clm}{\newtheorem{clm}{Clm}}{}
\@ifundefined{lem}{\newtheorem{lem}{Lem}}{}
\@ifundefined{rem}{\newtheorem{rem}{Rem}}{}
\@ifundefined{thm}{\newtheorem{thm}{Thm}}{}
\makeatother
\DeclareMathOperator{\ar}{ar}
\DeclareMathOperator{\supp}{supp}
\newcommand{\Add}[1]{#1|_\Gamma}
\newcommand{\Exists}[1]{\exists\,#1\ }
\newcommand{\Forall}[1]{\forall\,#1\ }
\newcommand{\GP}[1]{\left\langle#1\right\rangle_\Gamma}
\newcommand{\N}{\mathbb N}
\newcommand{\Prob}[1]{\Pr[#1]}
\newcommand{\acl}[1]{\lvert#1\rvert}
\newcommand{\alg}[1]{\langle#1\rangle}
\newcommand{\bb}{\mathsf}
\newcommand{\gp}[1]{\langle#1\rangle_\Gamma}
\title[Polynomial-Time Algorithms for Black-Box Distributive Expand]{Polynomial-Time Algorithms for Black-Box Distributive Expanded Groups}
\author{Mikhail Anokhin}
\date{Source: arXiv:2505.20497v2 (CC BY 4.0)}
\begin{document}
\maketitle

\noindent\emph{Source and licence.} Polynomial-Time Algorithms for Black-Box Distributive Expanded Groups. Version arXiv:2505.20497v2 (CC BY 4.0), distributed under the Creative Commons Attribution 4.0 International licence, as recorded in the arXiv licence metadata for this exact version and shown on the abstract page https://arxiv.org/abs/2505.20497v2. Full rights notice: licenses/arxiv-2505.20497-CC-BY-4.0.txt.\par\smallskip

\noindent\textit{Editorial scope.} Counted unit: the Thm whose source label is \texttt{t:compgensysofaddgr} in the article (source0.tex lines 1086--1096, source label \texttt{t:compgensysofaddgr}), delivered together with the article's own complete original proof of it (source0.tex lines 1097--1172, one \texttt{proof} environment of 76 lines). No mathematics is written, repaired or re-derived: statement and proof are the article's own text line for line. Required context retained uncounted: e:tau (lines 681--686); l:tauSeqtaugpS (lines 693--697); l:uniontauiSeqalgSSigma (lines 719--723); r:lenofchofsubgr (lines 863--869); l:Prrandsubprdonotgen (lines 891--902). Every definition, display and label of those lines is retained unchanged. Omitted from this package and not redistributed: 1--680, 687--692, 698--718, 724--862, 870--890, 903--1085, 1173--1370. Nothing inside a retained excerpt is omitted or altered: every retained body is delivered exactly as the article's lines are written, so no omission receipt of any kind accompanies this package. Citations: the retained excerpts carry no citation command at all, so no bibliography entry of the article is transcribed and no reference is redistributed. Reference adaptation: none was needed and none was made; every reference inside the delivered unit resolves inside it, so no unresolved reference or citation is produced. Printed slips kept literal: no printed word, symbol, hypothesis or number of the counted statement or of its proof was altered. Layout and class: the article's own class is \texttt{jgcc} with packages including lastpage, orcidlink, hyperref; the wrapper is the lane's pinned \texttt{amsart} editorial wrapper at 10pt on A4, which restates the theorem-like environments the retained text needs and adds the article's own macro definitions that the retained text needs (\texttt{\textbackslash Add}, \texttt{\textbackslash Exists}, \texttt{\textbackslash Forall}, \texttt{\textbackslash GP}, \texttt{\textbackslash N}, \texttt{\textbackslash Prob}, \texttt{\textbackslash acl}, \texttt{\textbackslash alg}, \texttt{\textbackslash ar}, \texttt{\textbackslash bb}, \texttt{\textbackslash gp}, \texttt{\textbackslash supp}), with extra wrapper packages (bm, bbm, dsfont, xspace, cancel, mathtools). The delivered numbering is the unit's own. Assets: no retained span contains a figure environment, a graphics-inclusion command, a PDF-page inclusion, a table or a TikZ picture; no image file is copied into this package, the package declares no asset dependency and no image file is redistributed with it. Licence: version arXiv:2505.20497v2 of this article is distributed under the Creative Commons Attribution 4.0 International licence, as recorded in the arXiv licence metadata for this exact version and shown on the abstract page https://arxiv.org/abs/2505.20497v2. A full rights notice is delivered at licenses/arxiv-2505.20497-CC-BY-4.0.txt. Characters: every retained excerpt is pure ASCII, so the delivered engine, XeLaTeX with its default Latin Modern text font, reports no missing character.
\begin{equation}
\label{e:tau}
\tau(S)
=\GP{S\cup\bigcup_{\omega\in\Omega}\omega(S^{\ar\omega})}
\quad\text{for any set }S\subseteq H.
\end{equation}

\begin{lem}
\label{l:tauSeqtaugpS}
Assume that $H$ is distributive. Then $\tau(S)=\tau(\gp S)$
for every set $S\subseteq H$.
\end{lem}

\begin{lem}
\label{l:uniontauiSeqalgSSigma}
For any set $S\subseteq H$, we have
$\bigcup_{i=0}^\infty\tau^i(S)=\alg S_\Sigma$.
\end{lem}

\begin{rem}
\label{r:lenofchofsubgr}
It is well known that if $G$ is finite, then any chain of
subgroups of $G$ has length at most~$\log_2\acl G$. This
follows from the fact that if $K$ is a proper subgroup of a
finite group $L$, then $\acl L\ge2\acl K$.
\end{rem}

\begin{lem}
\label{l:Prrandsubprdonotgen}
Assume that the group $G$ is finite and that any chain of
subgroups of $G$ has length at most $l\in\N$. Suppose $s\in
G^*$ is a generating system of~$G$. Also, let
$k\in\N\setminus\{0,1,2\}$ and let $r_1,\dots,r_{kl}$ be
independent random subproducts of~$s$. Then
\[
\Prob{\{r_1,\dots,r_{kl}\}\text{ does not generate }G}\le
e^{-(1-2/k)^2kl/4}.
\]
\end{lem}

\begin{thm}
\label{t:compgensysofaddgr}
Suppose $B$ is the algorithm described just before this
theorem. Let $\bb H=(H,\rho)$, $n$, $s$, and $O$ be as in
the description of~$B$. Then $\supp B^O(1^n,s)\subseteq H^*$
and
\begin{equation}
\label{e:PrgpBO1nseqH}
\Prob{\gp{B^O(1^n,s)}=H}\ge1-\frac n{c^n}.
\end{equation}
\end{thm}

\begin{proof}
We will use the notation introduced in the description of
the algorithm~$B$. It is evident that $\supp
B^O(1^n,s)\subseteq H^*$. Since $\acl H\le2^n$,
Remark~\ref{r:lenofchofsubgr} implies that any chain of
subgroups of $\Add H$ has length at most~$n$.

\begin{clm}
\label{cl:Prforalli}
We have
\begin{equation}
\label{e:Prforalli}
\Prob{\Forall{i\in\{1,\dots,n\}}(\gp{R_i}=\gp{s_{i-1}})}
\ge1-\frac n{c^n}.
\end{equation}
\end{clm}
\begin{proof}
We have already seen that any chain of subgroups of $\Add H$
has length at most~$n$. Therefore, by
Lemma~\ref{l:Prrandsubprdonotgen}, we have
$\Prob{\gp{R_i}\ne\gp{s_{i-1}}}\le e^{-(1-2/k)^2kn/4}\le
c^{-n}$ for every $i\in\{1,\dots,n\}$. Hence,
\[
\Prob{\Exists{i\in\{1,\dots,n\}}(\gp{R_i}\ne\gp{s_{i-1}})}
\le\frac n{c^n},
\]
which is equivalent to inequality~\eqref{e:Prforalli}.
\end{proof}

In the rest of the proof, we use the function $\tau$ from
the set of all subsets of $H$ to itself defined
by~\eqref{e:tau}.

\begin{clm}
Assume that $\gp{R_i}=\gp{s_{i-1}}$ for all
$i\in\{1,\dots,n\}$. Then $\gp{s_j}=\tau^j(\gp s)$ for every
$j\in\{0,\dots,n\}$.
\end{clm}
\begin{proof}
We proceed by induction on~$j$. Since $s_0=s$, we have
$\gp{s_0}=\gp s=\tau^0(\gp s)$. If $j\in\{1,\dots,n\}$ and
$\gp{s_{j-1}}=\tau^{j-1}(\gp s)$, then
\[
\gp{s_j}=\tau(R_j)=\tau(\gp{R_j})=\tau(\gp{s_{j-1}})
=\tau(\tau^{j-1}(\gp s))=\tau^j(\gp s),
\]
where the second equality follows from
Lemma~\ref{l:tauSeqtaugpS}. Thus, the claim holds.
\end{proof}

\begin{clm}
\label{cl:taungpseqH}
We have $\tau^n(\gp s)=H$.
\end{clm}
\begin{proof}
It is evident that
\[
\gp s=\tau^0(\gp s)\subseteq\tau^1(\gp
s)\subseteq\dots\subseteq H.
\]
Let $l$ be the smallest nonnegative integer such that
$\tau^l(\gp s)=\tau^{l+1}(\gp s)$ (it exists because $H$ is
finite). Then by induction on $i$, $\tau^i(\gp s)=\tau^l(\gp
s)$ for all $i\ge l+1$. By
Lemma~\ref{l:uniontauiSeqalgSSigma}, $\tau^l(\gp
s)=\tau^{l+1}(\gp s)=\dots=\alg{\gp s}_\Sigma=H$. Since any
chain of subgroups of $\Add H$ has length at most $n$, we
have $l\le n$. Thus, $\tau^n(\gp s)=H$.
\end{proof}

Inequality~\eqref{e:PrgpBO1nseqH} follows immediately from
Claims~\ref{cl:Prforalli}--\ref{cl:taungpseqH}. Indeed, if
$\gp{R_i}=\gp{s_{i-1}}$ for all $i\in\{1,\dots,n\}$, which
holds with probability at least $1-n/c^n$, then
$\gp{B^O(1^n,s)}=\gp{s_n}=\tau^n(\gp s)=H$.
\end{proof}

\end{document}
